% !PRINTED 752
\weberpage{752}

%<<CONT
have gradually been so completely removed that the proof can now be carried out by wholly elementary means and in the simplest way.

If, as already before, we denote by $\Pi(n)$ the product
\[
\Pi(n) = 1.2.3\dots n
\]
then, if $x$ is an arbitrary quantity,
\weq{1}{\frac{x^n}{\Pi(n)}}
will approach the limit zero as $n$ increases indefinitely, and indeed more strongly than the terms of a decreasing geometric series.

For if $k$ is an integer greater than the absolute value of $x$, then the absolute value of $x : h$ is a proper fraction whenever $h$ is equal to or greater than $k$. Consequently
\[
\left|\frac{x^n}{\Pi(n)}\right| = \left|\frac{x^k}{\Pi(k)}\, \frac{x}{k+1}\, \frac{x}{k+2} \cdots \frac{x}{n}\right| < \left|\frac{x^k}{\Pi(k)}\right| \left|\frac{x}{k}\right|^{n-k}.
\]
% [sic] -- the strict inequality tacitly assumes $n>k$ and $x\ne0$; $x=0$ is the trivial equality and convergence case.
From this it follows that the infinite series
\weq{2}{e^x = 1 + x + \frac{x^2}{\Pi(2)} + \frac{x^3}{\Pi(3)} + \cdots}
converges for all values of $x$, and by it we define the exponential function $e^x$, whose simplest properties we here assume from analysis. In particular, for any two values $x, y$ the relation
\[
e^{x+y} = e^x e^y,
\]
holds, from which it then follows that, for a positive integer $n$, $e^n$ is the $n$th power of the number
\weq{3}{e = 2 + \frac{1}{\Pi(2)} + \frac{1}{\Pi(3)} + \frac{1}{\Pi(4)} + \cdots = 2{,}718281828459\cdots\,.}

If now $r$ denotes any positive integer, we can represent formula (2) in the following way:
\weq{4}{\Pi(r)\, e^x = \Pi(r) + \frac{\Pi(r)}{\Pi(1)} x + \frac{\Pi(r)}{\Pi(2)} x^2 + \cdots + x^r + x^r\, U_r,}

\noindent\rule{4cm}{0.4pt}

Zahl\grqq{}. Sitzungsbericht der Berliner Akademie, 3. December 1885. The papers by \gesperrt{Hilbert}, \gesperrt{Hurwitz} and \gesperrt{Gordan} are all three to be found in volume 43 of Mathematische Annalen (1893), the first two also in the G\"ottinger Nachrichten of 1893.
%CONT>>
